%% ============================================================
%%  生成AI利用申告書 — テンプレート
%%  ハッカソン1 提出課題用 (塩澤 / 2026)
%%  コンパイル: docker run --rm -v "$(pwd):/workspace" -w /workspace \
%%              ghcr.io/paperist/texlive-ja:debian latexmk main.tex
%% ============================================================
\documentclass[a4paper,11pt,dvipdfmx]{jsarticle}

\usepackage[top=25mm,bottom=25mm,left=25mm,right=25mm]{geometry}
\usepackage{amssymb}
\usepackage{listings}
\usepackage{plistings}  % 日本語コメント混在時の文字幅補正
\usepackage[dvipdfmx,hidelinks]{hyperref}
\usepackage{pxjahyper}

\lstset{
  basicstyle=\ttfamily\small,
  frame=single,
  breaklines=true,
  xleftmargin=1em,
  aboveskip=6pt, belowskip=6pt,
  columns=flexible,
}

\newcommand{\cboxon}{$\boxtimes$}
\newcommand{\cboxoff}{$\square$}

\begin{document}

{\LARGE\bfseries 生成AI利用申告書}
\vspace{1em}

\begin{itemize}
  \item 氏名：
  \item 学籍番号：
  \item プログラムファイル名：
  \item 使用したAIツール（例：ChatGPT など）：
\end{itemize}

\section*{1. 何に使ったか}
今回の課題で、生成AIをどの部分に使いましたか。

（例：エラー修正／コード作成／アイデア出し など）

\section*{2. AIへの指示（プロンプト）}
AIにどのような質問・指示を出しましたか。重要なものを1〜2個書いてください。

（ここに記入）

\section*{3. AIの出力の使い方}
AIの出力をどのように使いましたか。
\begin{itemize}
  \item \cboxoff そのまま使った
  \item \cboxoff 一部修正して使った
  \item \cboxoff 参考にしただけ
\end{itemize}
$\rightarrow$ 上で選んだ理由を書いてください：

\section*{4. 正しいと確認した方法}
AIの出力が正しいと、どのように確認しましたか。

（例）
\begin{itemize}
  \item 実行して動作を確認した
  \item 自分で別の入力でテストした
  \item コードを読んで意味を理解した
\end{itemize}
※ 具体的に書くこと

\section*{5. 問題や間違い}
AIの出力に問題や間違いはありましたか。
\begin{itemize}
  \item \cboxoff あった
  \item \cboxoff なかった（※本当に？と思って確認した内容を書くこと）
\end{itemize}
確認した内容や気づいたこと：

\section*{6. 自分で工夫した点}
最終的なプログラムの中で、自分で考えて工夫したところを書いてください。

\section*{参考文献}

\end{document}
